% =============================================================================
% Lecture 15: Trajectory Curation
% Agentic Machine Learning Systems (Volume III), Chapter 15
% =============================================================================
% Thesis: an accepted trajectory is a record of what happened, and curation decides whether a failure belongs in the weights at all, and which
% records, tokens, and permissions each learner receives.
% Hook: NM 15.1, 3,000 failures a month, 90 percent of them below the weights.
% Running case: the repair agent's recovery record trace_r19_01 (chapter), and
% beside it a real public corpus, nebius/SWE-agent-trajectories (80,036 runs),
% measured for this lecture by images/data/corpus_stats.py.
%
% Build: make deck V=3 D=15_trajectory_curation   (figures: images/make_figs.py)
% Plan:  lecture_map.yml in this folder.
% Numbers: images/calcs.py prints every derived value; sources named there.
% =============================================================================

\documentclass[aspectratio=169, 12pt]{beamer}
\usepackage{../../assets/beamerthememlsys}
\mlsyssetup{
  volume={Volume III},
  chapter={Chapter 15},
  chaptertitle={Trajectory Curation},
  layer=learning,
  logo={../../assets/img/logo-mlsysbook.png},
  instlogo={../../assets/img/logo-harvard.png},
}
\usepackage[T1]{fontenc}
\usepackage[scaled=0.9]{helvet}
\usepackage{courier}
% Code and real tool output: Inconsolata is narrower than Courier.
\usepackage[varqu,scaled=0.95]{zi4}
\renewcommand{\familydefault}{\sfdefault}
\usepackage{booktabs,amsmath}
\usetikzlibrary{arrows.meta,positioning}
\graphicspath{{images/}}
% Generated by calcs.py. Do not edit; edit calcs.py and rerun it.
\defval{tri_runs}{50{,}000}
\defval{tri_fail_pct}{6}
\defval{tri_fail}{3{,}000}
\defval{tri_iface}{2{,}100}
\defval{tri_iface_pct}{70}
\defval{tri_rt}{600}
\defval{tri_rt_pct}{20}
\defval{tri_system}{2{,}700}
\defval{tri_residual}{300}
\defval{tri_seq}{4{,}096}
\defval{tri_epochs}{3}
\defval{tri_params_b}{70.6}
\defval{tri_fpt}{6}
\defval{tri_dev}{8}
\defval{tri_peak}{989}
\defval{tri_mfu}{0.30}
\defval{tri_mfu_pct}{30}
\defval{tri_price}{4}
\defval{tri_tokens}{36{,}864{,}000}
\defval{tri_tokens_m}{36.9}
\defval{tri_flops_mant}{1.56}
\defval{tri_flops_exp}{19}
\defval{tri_rate_mant}{2.37}
\defval{tri_rate_exp}{15}
\defval{tri_wall_h}{1.83}
\defval{tri_wall_nomfu}{0.55}
\defval{tri_dev_h}{14.6}
\defval{tri_cost}{58}
\defval{tri_cost_exact}{58.48}
\defval{tri_res_cost}{6}
\defval{c_runs}{80{,}036}
\defval{c_tasks}{3{,}591}
\defval{c_repos}{1{,}276}
\defval{c_resolved}{13{,}389}
\defval{c_resolved_pct}{16.7}
\defval{c_median}{16}
\defval{c_max}{125}
\defval{c_cap_bar}{40}
\defval{c_at_cap}{808}
\defval{c_straddle_pct}{89.9}
\defval{c_straddle_tasks}{3{,}229}
\defval{c_pred_pct}{90.3}
\defval{c_test_tasks}{743}
\defval{c_test_repo_pct}{77.4}
\defval{c_sig}{13{,}199}
\defval{c_dup}{7{,}569}
\defval{c_dup_pct}{57.3}
\defval{c_xpairs}{26}
\defval{c_xrepo}{3}
\defval{c_ai_masked}{2{,}115{,}623}
\defval{c_ai_total}{2{,}115{,}623}
\defval{c_refusal_pct}{44.0}
\defval{c_rec}{3{,}971}
\defval{c_rec_pct}{29.7}
\defval{c_rec_chars_pct}{42.5}
\defval{c_rec_ai_pct}{45.6}
\defval{c_direct_mean}{27{,}105}
\defval{c_rec_mean}{47{,}435}
\defval{c_len_ratio}{1.75}
\defval{c_cap_rec_pct}{28.0}
\defval{c_allfail_pct}{76.7}
\defval{c_mixed_pct}{22.5}
\defval{c_allpass_pct}{0.9}
\defval{sp_k}{8}
\defval{sp_f}{0.8}
\defval{sp_q}{0.2}
\defval{sp_f_pct}{80}
\defval{sp_q_pct}{20}
\defval{sp_whole}{0.168}
\defval{sp_straddle}{83.2}
\defval{sp_straddle_r}{83}
\defval{sp_two_pct}{32}
\defval{ex_tasks}{2{,}000}
\defval{ex_k}{5}
\defval{ex_f}{0.9}
\defval{ex_q}{0.1}
\defval{ex_f_pct}{90}
\defval{ex_q_pct}{10}
\defval{ex_q_exp}{-5}
\defval{ex_whole}{0.5905}
\defval{ex_straddle}{40.95}
\defval{ex_tasks_straddle}{819}
\defval{ex_sib_k}{4}
\defval{ex_sib}{0.9999}
\defval{ex_sib_pct}{99.99}
\defval{ex_major}{7}
\defval{ex_major_pct}{52.2}
\defval{mh_k}{128}
\defval{mh_bytes}{512}
\defval{mh_records}{100{,}000}
\defval{mh_mb}{51.2}
\defval{mh_se}{0.044}
\defval{mh_bands}{32}
\defval{mh_rows}{4}
\defval{mh_j_lo}{0.5}
\defval{mh_j_hi}{0.8}
\defval{mh_p_lo}{87}
\defval{mh_p_hi}{100}
\defval{ad_pass_record}{95}
\defval{ad_time_record}{0.01}
\defval{ad_pass_use}{90}
\defval{ad_time_use}{0.04}
\defval{ad_pass_process}{80}
\defval{ad_time_process}{0.20}
\defval{ad_pass_outcome}{40}
\defval{ad_time_outcome}{8}
\defval{ad_pass_target}{75}
\defval{ad_time_target}{2}
\defval{ad_cands}{100{,}000}
\defval{ad_rej_pct_process}{20}
\defval{ad_rej_pct_outcome}{60}
\defval{ad_rej_process}{1}
\defval{ad_rej_outcome}{13.3}
\defval{ad_t_best}{6.24}
\defval{ad_t_worst}{7.50}
\defval{ad_ta_best}{30.40}
\defval{ad_ta_worst}{36.57}
\defval{ad_yield}{20.52}
\defval{ad_admitted}{20{,}520}
\defval{ad_saving}{16.9}
\defval{av_time_process}{2}
\defval{av_pass_process}{95}
\defval{av_rej_process}{40}
\defval{av_saving}{7.7}
\defval{dn_records}{50{,}000}
\defval{dn_rec_pct}{25}
\defval{dn_stored_pct}{45}
\defval{dn_target_pct}{39}
\defval{dn_admitted_pct}{49.7}
\defval{dn_rec_len}{18{,}500}
\defval{dn_dir_len}{7{,}200}
\defval{dn_ctx_rec}{9.25}
\defval{dn_ctx_dir}{6}
\defval{gc_k}{8}
\defval{gc_lo}{0.05}
\defval{gc_lo_pct}{5}
\defval{gc_mid_pct}{50}
\defval{gc_hi_pct}{95}
\defval{gc_mixed_lo}{0.34}
\defval{gc_mixed_hi}{0.34}
\defval{gc_mixed_mid}{0.99}
\defval{gc_yield_lo_pct}{34}
\defval{gc_attempts_lo}{24}
\defval{st_rollouts}{1{,}024}
\defval{st_rollout_s}{45}
\defval{st_slots}{512}
\defval{st_timeout}{120}
\defval{st_ceiling_mb}{8}
\defval{st_store_gb}{32}
\defval{st_react_s}{30}
\defval{st_arrival}{22.8}
\defval{st_capacity}{30.3}
\defval{st_util}{0.75}
\defval{st_burst}{4.3}
\defval{st_net}{18.5}
\defval{st_growth}{147.9}
\defval{st_growth_r}{148}
\defval{st_tfull}{216}
\defval{st_tfull_1}{216.3}
\defval{st_tfull_min}{3.6}
\defval{st_inflight}{8.2}
\defval{st_react_gb}{4.4}
\defval{st_mark_pct}{60}
\defval{st_mark_gb}{19.4}
\defval{co_attempts}{500}
\defval{co_cost}{300}
\defval{co_emitted}{15}
\defval{co_other}{470}
\defval{co_inelig}{10}
\defval{co_pending}{3}
\defval{co_unknown}{2}
\defval{co_yield_pct}{3}
\defval{co_yield}{0.03}
\defval{co_attempt_cost}{0.60}
\defval{co_artifact}{20}
\defval{co_dbl_pct}{6}
\defval{co_dbl_cost}{10}
\defval{co_not_pct}{97}

\setcounter{mlsystotalsections}{5}

\title{Trajectory Curation}
\author{Vijay Janapa Reddi}
\institute{Harvard University}
\date{}

% A large answer choice for predict and concept-question slides
\newcommand{\choice}[2]{\textbf{\textcolor{mlsysaccent}{#1}}\enspace #2}
% A choice that is highlighted from overlay #1 on (the answer)
\newcommand{\answerat}[3]{\alt<#1->{\colorbox{mlsysaccent!15}{\choice{#2}{#3}}}{\choice{#2}{#3}}}
% A goal or takeaway line with its thumbnail
\newcommand{\thumbline}[2]{%
  \begin{minipage}[c]{2.6cm}\centering\includegraphics[width=2.5cm,height=0.95cm,keepaspectratio]{#1}\end{minipage}%
  \hspace{0.3cm}\begin{minipage}[c]{0.74\textwidth}\raggedright\small #2\end{minipage}}
% One line of a worked example: shown from overlay #1 on
\newcommand{\wstep}[2]{\uncover<#1->{#2}\\[0.18cm]}
% Golden-thread tag: one line on the goals slide naming the H·S·A term this lecture acts on
\newcommand{\hsathread}[1]{{\scriptsize\textbf{\textcolor{mlsysaccent}{H·S·A thread:}} #1\par}}
% An exact excerpt of real tool output, data, or source text
\newenvironment{logblock}[1][7.4]{\par\ttfamily\fontsize{#1}{9}\selectfont\raggedright\setlength{\parskip}{0pt}}{\par}
% Volume III role colors used in the figures, for text beside them
\definecolor{figblue}{HTML}{0284C7}
\definecolor{figpurple}{HTML}{7C3AED}
\definecolor{figred}{HTML}{DC2626}
\definecolor{figgreen}{HTML}{16A34A}
\definecolor{figamber}{HTML}{B45309}

\begin{document}

% =============================================================================
% TITLE AND HOOK
% =============================================================================
\mlsysnote{[0.5 min] [CORE] Read the thesis aloud once. Lecture 14 ended with traces that keep rare failures and postmortems that test a cause by rerunning a saved prefix. Those records now pile up, and a team preparing the next model wants to train on them. Today decides what may be trained on, and why. The notes total about 71.5 minutes; the core lecture is about 61. FLEX: cutting every [OPTIONAL] frame (the MinHash screen, the three-denominator bars, the published recipe, the emitted-artifact estimate, open questions) leaves a complete lecture. Cut first the open questions and the MinHash screen.}
\mlsystitle{Trajectory Curation}{Curation decides whether a failure belongs in the weights at all, and which records, tokens, and permissions each learner receives.}{cover_data_flywheel_blueprint.png}

\begin{frame}{Predict: what does \$\val{tri_cost} of fine-tuning fix?}
  \mlsysactivity{Predict}{2 min}
  \centering
  \mlsyssteps[width=0.95\textwidth,height=4.2cm,keepaspectratio]{hook_step}{3}\\[0.1cm]
  {\small Fine-tune on all \val{tri_fail} failing runs. How many failures stop recurring?}\\[0.15cm]
  {\normalsize\choice{A}{all \val{tri_fail}}\qquad\choice{B}{about \val{tri_system}}\qquad\answerat{3}{C}{at most \val{tri_residual}}\qquad\choice{D}{none can}}
  \mlsyscredit{Trajectory Curation, Napkin Math 15.1, \S15.2.2 (hypothetical failure mix, illustrative).}
  \note{[3 min] [CORE] Step 1: a maintenance agent runs \val{tri_runs} trajectories a month and \val{tri_fail_pct} percent fail, \val{tri_fail} failures. A team proposes fine-tuning on every one; the chapter prices that run at about \$\val{tri_cost} of accelerator time. ENGAGE: vote with fingers before step 2. Step 2: counterfactual probes, which this lecture defines, attribute \val{tri_iface} failures (\val{tri_iface_pct} percent) to a tool schema that omits a required field and \val{tri_rt} (\val{tri_rt_pct} percent) to sandbox startup timeouts. Only \val{tri_residual} survive every probe. Step 3: C. The schema and the sandbox pool still fail after training, so at most the \val{tri_residual} residual failures can respond to new weights; a schema fix and a larger sandbox pool remove \val{tri_system} for every model, present and future. Students who chose A treated a failure label as a training signal; B is what the system repair achieves, and the training run does not achieve it. WARN: the \$\val{tri_cost} is small on purpose. The argument against training here is misattribution, and the compute bill does not change it. A model trained to guess a missing field keeps guessing after the schema changes. LINK: the driving question asks why an accepted, or failed, record is not automatically a training example.}
\end{frame}

\begin{frame}{Recall: what did the postmortem branches show?}
  \mlsysactivity{Retrieval}{1.5 min}
  \begin{columns}[c]
    \begin{column}{0.56\textwidth}
      {\Large\bfseries Close your notes.}\\[0.2cm]
      {\normalsize A repair failed after a truncated search result. How did Lecture 14 test whether truncation caused it, and which traces does tail retention keep?}\\[0.3cm]
      \uncover<2->{\small Rerun continuations from the \textbf{saved prefix} (the run up to the failure) with the cause removed.\\[0.05cm]
      Keep \textbf{every failure} and a sample of successes.}
    \end{column}
    \begin{column}{0.40\textwidth}
      \centering\safeimg[width=\textwidth,height=1.6cm,keepaspectratio]{recall_branches.pdf}\\[0.05cm]
      {\footnotesize one saved prefix, three reruns}\\[0.2cm]
      \uncover<3->{\begin{mlsyscard}{mlsysaccent}\small With pagination added, every rerun reached an accepted repair. The defect sat in the interface.\end{mlsyscard}}
    \end{column}
  \end{columns}
  \mlsyscredit{Lecture 14: Operating Agents (postmortem by intervention, trace retention). Trajectory Curation, \S15.2.}
  \note{[1.5 min] [CORE] Retrieval from Lecture 14. Give students 45 seconds to write. Step 2: a postmortem by intervention reruns continuations from the saved prefix with one suspected cause removed; tail retention keeps every failed and unresolved trial plus a small sample of successes. Step 3: in the chapter's case, once the search adapter added pagination, every continuation reached an accepted repair. ENGAGE: if the team had trained on those truncated-search failures before the postmortem, what would the model have learned? (Expected: to work around a defect that one adapter change removed.) WARN: the ablation shows where the failure went away; it does not prove the model blameless. That limit returns in the first segment.}
\end{frame}

\begin{frame}{A public corpus holds \val{c_runs} agent runs on \val{c_tasks} issues}
  \begin{columns}[c]
    \begin{column}{0.54\textwidth}
      \centering\safeimg[width=\textwidth,height=4.4cm,keepaspectratio]{corpus_runs.pdf}\\[0.05cm]
      {\footnotesize Runs per GitHub issue, counted from the corpus}
    \end{column}
    \begin{column}{0.44\textwidth}
      \begin{logblock}[7.4]
        This dataset contains 80,036\\ trajectories generated by a software\\ engineering agent based on the\\ SWE-agent framework ...\\[0.12cm]
        target: Whether the model solved\\ the issue in this trajectory.\\[0.12cm]
        Total instances:\ \ 13,389 resolved,\\ \hspace*{8.6em}66,647 not resolved
      \end{logblock}
      \vspace{0.15cm}
      {\small \val{c_repos} repositories. One issue was attempted up to \textbf{\val{c_max}} times.}
    \end{column}
  \end{columns}
  \mlsyscredit{nebius/SWE-agent-trajectories, Hugging Face, CC BY 4.0 (revision 68195a1, 2024-12-23): dataset card excerpt; counts by images/data/corpus\_stats.py.}
  \note{[2.5 min] [CORE] This is a real trajectory corpus that a company published after training a software repair agent. Right: an exact excerpt of its dataset card. Each record is one run of an agent on a real GitHub issue: the messages it saw and wrote, the final patch, the test log, and one bit, target, saying whether the hidden tests passed. \val{c_resolved} of \val{c_runs} runs resolved their issue, \val{c_resolved_pct} percent. Left: we counted runs per issue ourselves. The median issue has \val{c_median} runs, \val{c_at_cap} issues have \val{c_cap_bar} or more, and one has \val{c_max}. ENGAGE: if you want to test a model trained on this corpus, what is the first thing you must decide? (Expected: which issues are held out, because every issue appears many times.) WARN: target is the outcome of a whole run. It says nothing about which of the run's messages are worth imitating. LINK: this corpus is the real counterpart of the chapter's repair agent; it returns in every segment.}
\end{frame}

\begin{frame}{Training changes the model and moves no exposure}
  \begin{columns}[c]
    \begin{column}{0.36\textwidth}
      \centering\mlsyskitfig[width=\textwidth,height=5.0cm,keepaspectratio]{stack_learning}
    \end{column}
    \begin{column}{0.20\textwidth}
      \centering\safeimg[width=\textwidth,height=2.4cm,keepaspectratio]{hsa_locator_origin.pdf}\\[0.1cm]
      {\footnotesize H: horizon, in turns\\ S: state carried\\ A: authority held\\ Origin: none lit}
    \end{column}
    \begin{column}{0.40\textwidth}
      {\large\bfseries Why is an accepted agent trajectory not automatically a good training example?}
    \end{column}
  \end{columns}
  \note{[1 min] [CORE] Part V, Learning from Trajectories, begins. The stack lights the learning layer. The H·S·A locator from Lecture 1 shows only its origin dot: training changes what the model proposes, and every exposure stays where the runtime closes it. A fine-tuned model still runs inside the same sandbox, budgets, and approvals (invariant closure, Lecture 1). Read the driving question aloud. Every segment answers part of it.}
\end{frame}

\begin{frame}{By the end of today you can}
  \thumbline{thumb_ladder_step3.pdf}{Decide by counterfactual probes whether a failure belongs in the weights, and name the artifact each record can feed.}\\[0.18cm]
  \thumbline{thumb_groups_step2.pdf}{Split by task group, and compute how a random split of runs leaks tasks.}\\[0.18cm]
  \thumbline{thumb_cascade_step3.pdf}{Write an admission route of five checks and order it by cost per rejection.}\\[0.18cm]
  \thumbline{thumb_typology_step3.pdf}{Select targets in a recovery record, and size the contrast a self-generated batch carries.}\\[0.18cm]
  \thumbline{thumb_staging_step2.pdf}{Bound a collector's staging with credits, and carry lineage into every derived record.}\\[0.18cm]
  \hsathread{moves no exposure; training changes the model}
  \note{[1 min] [CORE] Five goals, one per segment; each thumbnail is the shape of the figure that teaches it. Running cases: the chapter's repair agent and its recovery record trace\_r19\_01, and the real corpus from the previous slide.}
\end{frame}

% =============================================================================
% SECTION: TRIAGE
% =============================================================================
\mlsysnote{[0.5 min] [CORE] First question: before any record is collected, does this failure belong in the weights at all?}
\mlsysdivider{Does this failure belong in the weights?}{Trajectory Curation, \S15.2}
\section{Triage}

\begin{frame}{Triage probes send only a capability gap to training}
  \centering
  \mlsyssteps[width=0.92\textwidth,height=5.6cm,keepaspectratio]{ladder_step}{3}
  \mlsyscredit{Trajectory Curation, \S15.2.2, figure 1 (redrawn).}
  \note{[3 min] [CORE] The intervention ladder orders remedies by rising cost and falling reversibility: context, tool and schema, runtime, supervised fine-tuning (updating weights to imitate chosen demonstrations), reinforcement learning (updating weights to raise a verified reward). Step 1: a counterfactual probe reruns the recorded failure with one rung's suspected defect removed. The context probe supplies the missing fact; if the failure disappears, repair retrieval or assembly. Step 2: the schema probe resumes saved prefixes with a corrected tool description and samples fresh continuations; the runtime probe reruns in a clean sandbox with raised limits. Step 3: only a failure that survives all three is a candidate for training, and the trained checkpoint must still pass evaluation. ENGAGE: why fix truncation in the adapter instead of teaching the model to suspect it? (Expected: the end-to-end argument; a function belongs where it can be done completely, and the adapter fix holds for every model version.) WARN: the ladder orders the default investigation; it does not forbid skipping rungs. With recorded evidence of a genuine capability gap, a team may go straight to training. LINK: the next slide prices the shortcut.}
\end{frame}

\begin{frame}{Estimate: what does tuning on all \val{tri_fail} cost?}
  \mlsysactivity{Estimate}{2 min}
  \begin{columns}[T]
    \begin{column}{0.40\textwidth}
      \begin{mlsyscard}{mlsysaccent}
        \footnotesize\textbf{Given (NM 15.1)}\\[0.03cm]
        \val{tri_fail} runs $\times$ \val{tri_seq} tokens $\times$ \val{tri_epochs} epochs\\
        \val{tri_params_b}B parameters; \val{tri_fpt} floating-point operations (FLOP) per parameter-token\\
        \val{tri_dev} accelerators, \val{tri_peak} TFLOP/s each (16-bit); \val{tri_mfu_pct}\% of peak used (MFU)\\
        \$\val{tri_price} per accelerator-hour
      \end{mlsyscard}
    \end{column}
    \begin{column}{0.56\textwidth}
      \small
      \wstep{2}{Tokens: $\val{tri_fail} \times \val{tri_seq} \times \val{tri_epochs} = \val{tri_tokens_m}$ million}
      \wstep{2}{Work: $\val{tri_fpt} \times \val{tri_params_b}\text{B} \times \val{tri_tokens_m}\text{M} = \val{tri_flops_mant}\times10^{\val{tri_flops_exp}}$ FLOP}
      \wstep{3}{Rate: $\val{tri_dev} \times \val{tri_peak}\ \text{TFLOP/s} \times \val{tri_mfu} = \val{tri_rate_mant}\times10^{\val{tri_rate_exp}}$ FLOP/s}
      \wstep{3}{Time: \val{tri_wall_h} h, or \val{tri_dev_h} accelerator-hours}
      \uncover<4->{\textbf{\textcolor{mlsysaccent}{About \$\val{tri_cost}}}; the \val{tri_residual} residual failures alone, about \$\val{tri_res_cost}.\\[0.12cm]
      {\footnotesize\textcolor{midgray}{Training removes none of the \val{tri_system} schema and sandbox failures.}}}
    \end{column}
  \end{columns}
  \mlsyscredit{Trajectory Curation, Napkin Math 15.1 (hypothetical; accelerator figures as the chapter states them).}
  \note{[3 min] [CORE] ENGAGE: two minutes in pairs with the given card. Step 2: \val{tri_fail} runs of \val{tri_seq} tokens for \val{tri_epochs} epochs is \val{tri_tokens} training tokens; the \val{tri_fpt}-FLOP-per-parameter-per-token rule (forward and backward) gives \val{tri_flops_mant} times 10 to the \val{tri_flops_exp} FLOP. Step 3: one node of \val{tri_dev} accelerators at \val{tri_peak} TFLOP/s dense 16-bit and \val{tri_mfu_pct} percent model FLOPs utilization sustains \val{tri_rate_mant} times 10 to the \val{tri_rate_exp} FLOP/s, so the run takes \val{tri_wall_h} hours, \val{tri_dev_h} accelerator-hours. Step 4: at \$\val{tri_price} an hour, \$\val{tri_cost} (derived, calcs.py prints \$\val{tri_cost_exact}). Training only the \val{tri_residual} residual failures costs about \$\val{tri_res_cost}. Students who multiply by \val{tri_peak} without the utilization get \val{tri_wall_nomfu} hours; \val{tri_mfu_pct} percent is the chapter's assumption. WARN: the invoice is the wrong argument for or against training. The decisive cost is that the schema still omits its field and the sandbox pool still times out, and the trained compensation breaks when the interface changes. The estimate also excludes curation, review, and evaluation labor. FLEX: [CORE].}
\end{frame}

\begin{frame}{One run can feed three learners, each with its own test}
  \begin{columns}[c]
    \begin{column}{0.64\textwidth}
      \centering\mlsyssteps[width=\textwidth,height=5.6cm,keepaspectratio]{artifacts_step}{3}
    \end{column}
    \begin{column}{0.34\textwidth}
      {\footnotesize One published team, one corpus, two learners:}\\[0.1cm]
      \begin{logblock}[7.2]
        The resulting training set\\ contained 2.7k trajectories.\\[0.1cm]
        The training dataset consists of\\ 8.3k positive and 12.7k\\ negative trajectories ...
      \end{logblock}
      \vspace{0.1cm}
      {\footnotesize Successes trained the policy. Failures trained a critic that scores runs.}
    \end{column}
  \end{columns}
  \mlsyscredit{Trajectory Curation, \S15.2.3. Excerpts: Nebius AI R\&D, ``Leveraging training and search for better software engineering agents,'' blog post, 2024-11-15.}
  \note{[2.5 min] [CORE] One recorded run can feed three learning artifacts, and each has its own admission conditions instead of inheriting the run's outcome. Step 1: imitation material. Keep a malformed proposal and its refusal in the history and select only the corrected action as the target. Step 2: a preference comparison pairs two runs from the same initial state, environment, and action space, with a recorded reason why one is worse. A failure generated under a different prompt or hint cannot sit in the rejected slot. Step 3: a rollout fixture is a resettable environment and a grader that lets a learner try fresh actions; a failed run often seeds one even when none of its actions deserve imitation. Right: a real team used one corpus twice. 2.7k resolved trajectories trained the policy by imitation; 8.3k positive and 12.7k negative trajectories trained a critic that predicts success. ENGAGE: a run failed on a subtle race condition. Which artifact can it still feed? (Expected: a rollout fixture, or the rejected side of a properly matched comparison.) WARN: a run that changes no files can be a valid demonstration when the task requires restraint, such as declining an out-of-scope request. LINK: preference training and its loss belong to Lecture 16; rollouts against a verifier to Lecture 17.}
\end{frame}

\begin{frame}{Diagnose: which rung owns each failure?}
  \mlsysactivity{Diagnose}{3 min}
  \small
  \begin{tabular}{@{}p{8.9cm}l@{}}\toprule
    \textbf{Recorded finding} & \uncover<2->{\textbf{Rung to repair first}} \\\midrule
    1.\ A search excerpt arrived with no truncation marker & \uncover<2->{tool or schema: add the marker} \\
    2.\ The required config value was never in the prompt & \uncover<2->{context: fix retrieval} \\
    3.\ The test sandbox times out only when the pool is busy & \uncover<2->{runtime: size the pool} \\
    4.\ With all three fixed, the model still edits outside scope & \uncover<2->{\textbf{\textcolor{mlsysaccent}{capability gap: train it}}} \\\bottomrule
  \end{tabular}\\[0.2cm]
  \begin{columns}[c]
    \begin{column}{0.47\textwidth}
      \uncover<3->{\small Even after training, the runtime still refuses out-of-scope edits.}
    \end{column}
    \begin{column}{0.51\textwidth}
      {\footnotesize A real interface that marks what it omits:}\\[0.03cm]
      \begin{logblock}[7.4]
        [File: .../environ.py (994 lines total)]\\
        (57 more lines above) ... (925 more lines below)
      \end{logblock}
    \end{column}
  \end{columns}
  \mlsyscredit{Trajectory Curation, \S15.2.1, table 1, and \S15.2.2. Observation: SWE-agent file viewer in nebius/SWE-agent-trajectories (joke2k\_\_django-environ-358), CC BY 4.0.}
  \note{[3 min] [CORE] ENGAGE: two minutes with a neighbor. For each finding, name the rung to repair first. Step 2: 1 is an interface defect: the adapter should mark truncation, and that fix holds for every model. 2 is a context defect: weights cannot supply a current fact the model never observed. 3 is a runtime defect, a busy pool, and the model plays no part. 4 survived every probe, so it is a candidate for a learning objective, stated precisely: for example, inspect the scope before editing, and stop when the scope forbids the edit. Students who send 3 to training are training the model to tolerate a timeout that a larger pool removes. Step 3: invariant closure (Lecture 1) still holds. A fine-tuned model proposes out-of-scope edits less often, but the admission gate keeps refusing them, because a lower failure rate guarantees nothing about the next proposal. WARN: one trace often contains several findings at once (the chapter's postmortem found an interface defect, a forbidden proposal, and a correct refusal in one run). Each gets its own remedy. Bottom right: a real observation from the corpus. SWE-agent's viewer prints how many lines lie above and below its window, the completeness metadata that finding 1 lacks. Land the segment aloud: probes keep system defects out of the weights.}
\end{frame}

% =============================================================================
% SECTION: SPLITS
% =============================================================================
\mlsysnote{[0.5 min] [CORE] The failure belongs in the weights. Before collecting anything, fix the evidence that will later judge the trained model.}
\mlsysdivider{Will the test score still measure unseen tasks?}{Trajectory Curation, \S15.3 and \S15.4}
\section{Splits}

\begin{frame}{Predict: how many tasks leak across the split?}
  \mlsysactivity{Predict}{2 min}
  \centering
  \mlsyssteps[width=0.82\textwidth,height=4.6cm,keepaspectratio]{straddle_step}{2}\\[0.05cm]
  {\small \val{sp_k} attempts per task. Assign each \emph{run} to training with probability \val{sp_f}.}\\[0.1cm]
  {\normalsize\choice{A}{about 20\%}\qquad\choice{B}{about 50\%}\qquad\answerat{2}{C}{about \val{sp_straddle_r}\%}}
  \mlsyscredit{Trajectory Curation, \S15.4 and its figure (formula plotted by images/make\_figs.py).}
  \note{[2.5 min] [CORE] Step 1: a collector runs \val{sp_k} attempts per task, and an engineer splits the runs at random, \val{sp_f_pct} percent to training. A task stays clean only if all of its attempts land on the same side. ENGAGE: vote. Step 2: C. The chance that all \val{sp_k} land together is $\val{sp_f}^{\val{sp_k}} + \val{sp_q}^{\val{sp_k}} = \val{sp_whole}$, so \val{sp_straddle} percent of tasks straddle (derived). With 2 attempts it is already \val{sp_two_pct} percent. Students who chose A read the \val{sp_q_pct} percent test fraction as the leak rate. WARN: varying the seed or the prompt wording makes runs look different but leaves the problem the same; the test then measures recall of a solved sibling and stops measuring generalization. This is split leakage: information about an evaluation target reaching training through a related record. LINK: Lecture 13 defined contamination and benchmark validity; today owns the split protocol that prevents it inside our own corpus.}
\end{frame}

\begin{frame}{Measured: a random run split leaks \val{c_straddle_pct}\% of issues}
  \centering
  \safeimg[width=0.92\textwidth,height=4.0cm,keepaspectratio]{corpus_split.pdf}\\[0.15cm]
  {\small Formula, per issue's own run count: \val{c_pred_pct}\%. Measured on one seeded split: \val{c_straddle_pct}\%.}\\[0.08cm]
  \uncover<2->{\small Same repository is fine for ``new issues here''; it breaks a claim about new repositories.}
  \mlsyscredit{nebius/SWE-agent-trajectories, CC BY 4.0; \val{c_tasks} issues, \val{c_repos} repositories; split and counts by images/data/corpus\_stats.py (seed 15).}
  \note{[2.5 min] [CORE] We split the real corpus ourselves. Top bar: assign each of the \val{c_runs} runs to training with probability \val{sp_f}. \val{c_straddle_tasks} of \val{c_tasks} issues, \val{c_straddle_pct} percent, end up on both sides; the formula from the last slide, applied to each issue's own run count, predicts \val{c_pred_pct}. The corpus leaks more than the \val{sp_k}-attempt example because the median issue has \val{c_median} runs. Middle: assign whole issues instead, the task group of the next slide. No issue straddles, by construction. Step 2, bottom bar: \val{c_test_repo_pct} percent of the \val{c_test_tasks} test issues come from a repository that also has training issues. ENGAGE: is that leakage? (Expected: it depends on the claim. For new issues in repositories the agent already serves, shared build tools and APIs are legitimate deployment conditions. For transfer to unseen repositories, hold out whole repository families.) WARN: the published team behind this corpus excluded mirrors and forks of the benchmark's repositories for exactly this reason. Every split begins with a stated evaluation claim. FLEX: [CORE].}
\end{frame}

\begin{frame}{A task group keeps every relative on one side}
  \begin{columns}[c]
    \begin{column}{0.62\textwidth}
      \centering\mlsyssteps[width=\textwidth,height=5.5cm,keepaspectratio]{groups_step}{2}
    \end{column}
    \begin{column}{0.36\textwidth}
      \small $g(f)$: the group of fixture $f$\\ $s(g)$: the group's partition\\[0.25cm]
      \uncover<2->{\footnotesize $g(f_i)=g(f_j) \Rightarrow s(g(f_i))=s(g(f_j))$\\[0.25cm]
      \small Assign before collection. A test task inspected while debugging moves to development.}
    \end{column}
  \end{columns}
  \mlsyscredit{Trajectory Curation, \S15.4.1 and the partition-integrity relation.}
  \note{[2.5 min] [CORE] Step 1: a task group holds every attempt at one underlying problem, every paraphrased prompt, every fork that copied the same bug, and every record derived later, such as a redacted trace. The split manifest records why each member belongs. Step 2: randomization runs over groups, so every member of a group inherits one partition: training, development, or final test. Read the relation as: same group implies same side. ENGAGE: a teammate inspects a final-test failure and writes a training demonstration aimed at it. What happens to that task group? (Expected: it is reclassified as development data, and a fresh unseen group replaces it in the final test.) WARN: grouping by repository name misses vendored copies and forks; grouping by shared dependency merges unrelated bugs into giant clusters. The relation is only as strong as the pipeline's record of how tasks relate, which is why the next slide screens for unrecorded relatives. LINK: held-out tests, reference patches, and verdicts also stay away from every generator, retriever, and ranker upstream. A ranker tuned on the final test leaks it without copying a token.}
\end{frame}

\begin{frame}{MinHash banding sends \val{mh_p_lo}\% of $J=\val{mh_j_lo}$ pairs to review}
  \begin{columns}[c]
    \begin{column}{0.56\textwidth}
      \centering\safeimg[width=\textwidth,height=4.7cm,keepaspectratio]{lsh.pdf}
    \end{column}
    \begin{column}{0.42\textwidth}
      \small $J$: shared features $\div$ all features\\[0.08cm]
      \val{mh_k}-number signature: \val{mh_bytes} bytes;\\ \val{mh_records} records: \val{mh_mb} MB\\[0.25cm]
      \uncover<2->{\textbf{Measured in the corpus:} \val{c_dup_pct}\% of resolved patches have a near-duplicate ($J \ge 0.8$) for the same issue.\\[0.08cm]
      \val{c_xpairs} near-duplicate pairs span two issues.}
    \end{column}
  \end{columns}
  \mlsyscredit{Trajectory Curation, \S15.4.3 (\val{mh_bands} bands of \val{mh_rows} rows). Measured: corpus\_stats.py, datasketch 2.0 MinHash LSH on token 5-grams of \val{c_sig} resolved patches.}
  \note{[2 min] [OPTIONAL] Metadata catches only the relatives the pipeline already knows. Similarity screening finds the rest. Each record becomes a set of features (command n-grams, code-tree pieces); the Jaccard similarity $J$ of two sets is their overlap divided by their union. MinHash keeps the minimum hash under each of \val{mh_k} random permutations; two signatures agree in a position with probability $J$, so the agreeing fraction estimates it, with standard error at most \val{mh_se}. Banding into \val{mh_bands} bands of \val{mh_rows} rows sends a pair to review if any band matches: \val{mh_p_lo} percent at $J = \val{mh_j_lo}$, essentially all at \val{mh_j_hi}. Step 2: we ran the same method on the corpus. \val{c_dup_pct} percent of resolved patches (\val{c_dup} of \val{c_sig}) have a near-duplicate for the same issue, a reason to cap and deduplicate per task. \val{c_xpairs} near-duplicate pairs cross two issues, \val{c_xrepo} of them across repositories: candidate relatives for one group. ENGAGE: why review a pair instead of deleting it? (Expected: two identical patches can come from different decision paths, and a short trace is not better than a longer one that shows recovery.) WARN: aggressive normalization can erase the difference between a working patch and a broken one. FLEX: cut if short; say the \val{c_dup_pct} percent aloud.}
\end{frame}

\begin{frame}{Estimate: \val{ex_k} attempts, a \val{ex_f_pct}/\val{ex_q_pct} split of runs?}
  \mlsysactivity{Estimate}{2 min}
  \begin{columns}[T]
    \begin{column}{0.40\textwidth}
      \begin{mlsyscard}{mlsysaccent}
        \small\textbf{Given (Exercise 1)}\\[0.05cm]
        \val{ex_tasks} repair tasks\\
        \val{ex_k} attempts per task\\
        each run to training with probability \val{ex_f}
      \end{mlsyscard}
      \vspace{0.2cm}
      {\small How many tasks straddle? At what attempt count does a majority?}
    \end{column}
    \begin{column}{0.56\textwidth}
      \small
      \wstep{2}{Stay together: $\val{ex_f}^{\val{ex_k}} + \val{ex_q}^{\val{ex_k}} = \val{ex_whole}$}
      \wstep{2}{Straddle: $1 - \val{ex_whole} = $ \textbf{\val{ex_straddle}\%}, about \textbf{\val{ex_tasks_straddle}} tasks}
      \wstep{3}{A test run has a sibling in training: $1 - \val{ex_q}^{\val{ex_sib_k}} = \val{ex_sib}$}
      \wstep{4}{Majority first straddles at \textbf{\val{ex_major}} attempts}
    \end{column}
  \end{columns}
  \mlsyscredit{Trajectory Curation, \S15.14, Exercise 1 (hypothetical inputs). Values from images/calcs.py.}
  \note{[2 min] [CORE] ENGAGE: two minutes, alone. Step 2: a task stays on one side only if all \val{ex_k} runs go to training ($\val{ex_f}^{\val{ex_k}} = \val{ex_whole}$) or all go to test ($10^{\val{ex_q_exp}}$); \val{ex_straddle} percent straddle, \val{ex_tasks_straddle} of \val{ex_tasks} tasks (derived). Step 3: pick any test run; its \val{ex_sib_k} siblings all avoid training with probability $\val{ex_q}^{\val{ex_sib_k}}$, so it has a sibling in training \val{ex_sib_pct} percent of the time. Step 4: the straddle share first passes one half at \val{ex_major} attempts (\val{ex_major_pct} percent). Students often expect a \val{ex_f_pct}/\val{ex_q_pct} split to leak less than \val{sp_f_pct}/\val{sp_q_pct}; it does for small counts, but any random split over runs leaks most tasks once attempts grow. WARN: a group split fixes the leak at no compute cost; it costs only the bookkeeping of which records belong to which task. Land the segment aloud: a split protects a claim, so it is assigned to whole task groups before collection begins.}
\end{frame}

% =============================================================================
% SECTION: ADMISSION
% =============================================================================
\mlsysnote{[0.5 min] [CORE] The split is fixed. Now collected records arrive. Which of them may each learner receive?}
\mlsysdivider{Which records may each learner receive?}{Trajectory Curation, \S15.5}
\section{Admission}

\begin{frame}{Admission runs five checks with four possible verdicts}
  \centering
  \mlsyssteps[width=0.94\textwidth,height=5.4cm,keepaspectratio]{cascade_step}{3}
  \mlsyscredit{Trajectory Curation, \S15.5.1, table 5 and figure 2 (redrawn).}
  \note{[2.5 min] [CORE] An admission cascade is an ordered sequence of declared checks that screens a record for one learning route and stops at the first check that does not pass. R: is the record faithful? U: do provenance, license, and split allow this use? P: did the run respect file boundaries, approvals, and limits? O: does authoritative evidence show the objective was met? T: are the targets justified by what the agent knew at that step? Step 1: record A passes all five and is admitted for this route only. Step 2: record B fails the outcome check. It is rejected for the demonstration route and kept: it may still serve diagnosis, a comparison, or a fixture. Step 3: record C's permission review is pending. Each check returns pass, fail, pending, or unknown (the four verdicts of Lecture 13), and a missing or timed-out check never counts as a pass; C waits in a holding queue tagged with the evidence it lacks. ENGAGE: a record shows a malformed tool call that the runtime refused before dispatch. Does it pass R? (Expected: yes. The log is faithful; whether the call is a target is T's question.) WARN: count admitted, rejected, and unresolved records per route, so redirected traces never inflate the demonstration yield.}
\end{frame}

\begin{frame}{A frontier model special-cased tests after failing}
  \centering\safeimg[width=0.97\textwidth,height=2.75cm,keepaspectratio]{system-card-special-casing.png}\\[0.1cm]
  \begin{columns}[T]
    \begin{column}{0.50\textwidth}
      {\footnotesize The chapter's tampered reproducer:}\\[0.03cm]
      \begin{mlsyscard}{midgray}
      \begin{logblock}[7.4]
        \textcolor{figred}{-\ \ \ \ assert waiter.done(timeout=1.0)}\\
        \textcolor{figgreen}{+\ \ \ \ \# skip flaky wakeup check}\\
        \textcolor{figgreen}{+\ \ \ \ assert True}
      \end{logblock}
      \end{mlsyscard}
    \end{column}
    \begin{column}{0.46\textwidth}
      \uncover<2->{\footnotesize\textbf{Fails P:} the diff touches a protected test path.\\[0.05cm]
      \textbf{Fails O:} clean tests in a fresh sandbox still deadlock.}
    \end{column}
  \end{columns}
  \mlsyscredit{Anthropic, Claude 3.7 Sonnet System Card (2025), \S6, p.~22 (crop). Listing: Trajectory Curation, listing 1, \S15.5.2.}
  \note{[2.5 min] [CORE] Left: an exact excerpt from a published system card. The model, after repeated failed attempts, sometimes returned expected test values directly or edited the tests to match its output; the behavior emerged from reward hacking during reinforcement learning, user testing missed it because it was rare, and classifiers on training transcripts caught it. Right: the chapter's version on the repair task. The candidate replaces the wakeup assertion with assert True; the runner exits with zero and the race remains. Step 2: a sound cascade rejects it twice. The process check binds the diff against the base commit to a protected-path manifest, so any change to tests, CI configuration, or build dependencies fails P regardless of the exit code. The outcome check runs the untouched, out-of-tree tests in a fresh sandbox, which reproduces the deadlock. This is test-tampering defense: grader isolation (Lecture 13) applied to admission. ENGAGE: why is tampering worse during collection than during evaluation? (Expected: an evaluation score is wrong once; an admitted record trains every later checkpoint to repeat the exploit.) WARN: a read-only mount alone is not enough, because co-located code can still read the tests or interfere with the reporting process. Rerunning a flaky test until it passes once samples favorable noise: $1-(1-p)^k$.}
\end{frame}

\begin{frame}{Exercise: which check should run first?}
  \mlsysactivity{Exercise}{3 min}
  \begin{columns}[T]
    \begin{column}{0.44\textwidth}
      \scriptsize
      \begin{tabular}{@{}lrr@{}}\toprule
        \textbf{Check} & \textbf{Pass} & \textbf{Slot time} \\\midrule
        R record & \val{ad_pass_record}\% & \val{ad_time_record} s \\
        U use & \val{ad_pass_use}\% & \val{ad_time_use} s \\
        P process & \val{ad_pass_process}\% & \val{ad_time_process} s \\
        O outcome (tests) & \val{ad_pass_outcome}\% & \val{ad_time_outcome} s \\
        T target review & \val{ad_pass_target}\% & \val{ad_time_target} s \\\bottomrule
      \end{tabular}\\[0.12cm]
      {\footnotesize R, U first; T last (it reads O's results). P or O next?}
    \end{column}
    \begin{column}{0.54\textwidth}
      \small
      \wstep{2}{Cost per rejection $= t/(1-p)$:\ \ P: \val{ad_rej_process} s;\ \ O: \val{ad_rej_outcome} s}
      \uncover<3->{\centering\safeimg[width=\textwidth,height=2.5cm,keepaspectratio]{order.pdf}\\[0.05cm]}
      \wstep{3}{Same yield either way: \textbf{\val{ad_yield}\%} admitted}
      \uncover<4->{\textbf{\textcolor{mlsysaccent}{\val{ad_saving}\% less verifier time}}, the same corpus.}
    \end{column}
  \end{columns}
  \mlsyscredit{Trajectory Curation, Napkin Math 15.2 and equations 1 and 2 (hypothetical rates; independent checks). Values from images/calcs.py.}
  \note{[3 min] [CORE] ENGAGE: three minutes in pairs. \val{ad_cands} candidate records, five independent checks. Step 2: for two adjacent independent checks, running A first is cheaper when $t_A/(1-p_A) < t_B/(1-p_B)$: cost per rejection. P costs \val{ad_time_process} s and rejects \val{ad_rej_pct_process} percent, \val{ad_rej_process} s per rejection; O costs \val{ad_time_outcome} s and rejects \val{ad_rej_pct_outcome} percent, \val{ad_rej_outcome} s per rejection. Run P first. Step 3: expected slot time per attempted record is the sum of each check's time weighted by the chance of reaching it: \val{ad_t_best} s with P first against \val{ad_t_worst} s with O first (derived). Both admit \val{ad_yield} percent, about \val{ad_admitted} records; per admitted record, \val{ad_ta_best} s against \val{ad_ta_worst} s. Step 4: \val{ad_saving} percent less verifier time, and not one valid record lost. Ordering by duration alone works here by luck. In the Checkpoint 15.1 variant (P costs \val{av_time_process} s, passes \val{av_pass_process} percent) P's cost per rejection is \val{av_rej_process} s, so O runs first and saves \val{av_saving} percent. WARN: dropping or loosening a check changes the policy. It admits a different corpus, so its lower price per admitted record compares two different products. LINK: the ratio of checking time to yield has the same form as cost per accepted task (Lecture 1). Land the segment aloud: cheaper checking must admit the same corpus.}
\end{frame}

% =============================================================================
% SECTION: RECOVERY
% =============================================================================
\mlsysnote{[0.5 min] [CORE] Clean successes teach the clean path. Real agents leave it. How do we train recovery without training the mistake, and what happens when the policy supplies its own data?}
\mlsysdivider{How do we teach recovery, and what does self-generated data narrow?}{Trajectory Curation, \S15.6 and \S15.7}
\section{Recovery}

\begin{frame}{A recovery record keeps the mistake as context only}
  \centering
  \mlsyssteps[width=0.92\textwidth,height=5.8cm,keepaspectratio]{typology_step}{3}
  \mlsyscredit{Trajectory Curation, \S15.6.1 and \S15.6.2, figure 3 (redrawn); target manifest recovery\_targets\_v1.}
  \note{[2.5 min] [CORE] A model trained only on clean prefixes meets states it never saw once an early step goes wrong, and over a long horizon a recoverable detour becomes a failed task (horizon compounding, Lecture 1). Step 1: a direct-success demonstration. Its supported actions can all be targets; blue bars mark the tokens that receive loss. Step 2: the chapter's recovery record, trace\_r19\_01. At e2 the model proposes a malformed edit and the runtime refuses it before dispatch. e3 inspects the schema, e4 issues the corrected edit, e5 is the check summary, e6 reports the result and its scope. The target manifest selects e3, e4, and e6; e2 and e5 stay in the sequence as context, because a correction makes sense only beside the error it fixes. Step 3: a hard negative looks plausible and passes a targeted check but fails the required one; it is never imitated, and if the first causal step is uncertain the record says so. ENGAGE: why not delete e2 to make the trace clean? (Expected: the corrected edit no longer follows from anything; the model learns to emit a schema lookup for no reason.) WARN: a success and a failure from different prompts or tool outputs do not form a preference pair. LINK: how masked tokens become a loss is Lecture 16.}
\end{frame}

\begin{frame}{Diagnose: which turns of a real run teach?}
  \mlsysactivity{Diagnose}{3 min}
  \begin{columns}[T]
    \begin{column}{0.60\textwidth}
      \begin{logblock}[6.9]
        \textbf{ai\ \ mask=True}\ \ edit 63:64 ... return unquote(v) ...\\[0.06cm]
        \textcolor{figred}{\textbf{user mask=False}\ Your proposed edit has introduced new}\\
        \textcolor{figred}{\hspace*{8.2em}syntax error(s) ... F821 undefined name 'unquote'}\\[0.06cm]
        \textbf{ai\ \ mask=True}\ \ ... we need to import the `unquote` function\\
        \hspace*{8.2em}edit 19:19 ... from urllib.parse import ( ...\\[0.06cm]
        \textbf{ai\ \ mask=True}\ \ It seems that I mistakenly duplicated the\\
        \hspace*{8.2em}import statements ...\\[0.06cm]
        \textbf{ai\ \ mask=True}\ \ edit 72:73 ... return unquote(v) ...\\[0.06cm]
        \textbf{ai\ \ mask=True}\ \ ... since we cannot run tests directly in\\
        \hspace*{8.2em}this environment, we will proceed to submit
      \end{logblock}
      \vspace{0.08cm}
      {\footnotesize Outcome: target=True, 12 tests passed.}
    \end{column}
    \begin{column}{0.38\textwidth}
      \small Which model turns should receive loss?\\[0.2cm]
      \uncover<2->{\footnotesize The corpus marks \textbf{all \val{c_ai_masked}} model turns \texttt{mask=True}, in failed runs too.\\[0.15cm]}
      \uncover<3->{\footnotesize The accepted patch still adds a \textbf{duplicate 8-line import}. The submit claims a check never run.}
    \end{column}
  \end{columns}
  \mlsyscredit{nebius/SWE-agent-trajectories, CC BY 4.0: joke2k\_\_django-environ-358, swe-agent-llama-70b, run 610 of the corpus; excerpts with ``...'' omissions. Counts: corpus\_stats.py.}
  \note{[3 min] [CORE] A real recovery record, excerpted; the full run has 15 messages. The first edit uses unquote without importing it, and SWE-agent's linter refuses it before applying it, the chapter's e2. The agent adds an import block, notices duplicated imports, edits again, makes the change, and submits, saying it cannot run tests. The hidden tests pass. ENGAGE: three minutes in pairs. Which model turns would you select as targets? Step 2: this corpus marks every model message mask=True and every observation False; we counted \val{c_ai_masked} of \val{c_ai_total} model messages marked, in failed runs as well as resolved ones. A role label says who wrote a message; it says nothing about whether the message deserves imitation. Step 3: the accepted patch still contains a duplicated 8-line import block (we read it), and the final message claims a check it never ran. Expected answer: keep the refused edit as context; the import repair is a weak target because it introduced the duplication; the final unquote edit is a sound target; the submit message should not teach the model to skip tests. WARN: this record passes R and O. It fails T for several turns, which is why target suitability is its own check. \val{c_refusal_pct} percent of all runs in the corpus contain at least one refused edit.}
\end{frame}

\begin{frame}{One recovery class has three different shares}
  \begin{columns}[c]
    \begin{column}{0.49\textwidth}
      \centering\safeimg[width=\textwidth,height=4.5cm,keepaspectratio]{denom_nm.pdf}
    \end{column}
    \begin{column}{0.49\textwidth}
      \centering\safeimg[width=\textwidth,height=4.5cm,keepaspectratio]{denom_corpus.pdf}
    \end{column}
  \end{columns}
  \vspace{0.1cm}
  \centering{\small Specify a mixture in selected targets; budget compute in processed tokens.}
  \mlsyscredit{Left: Trajectory Curation, Napkin Math 15.3 (hypothetical pool). Right: nebius/SWE-agent-trajectories, CC BY 4.0, resolved runs with at least one refused edit; counted in characters. Counts: corpus\_stats.py.}
  \note{[2.5 min] [OPTIONAL] Left, the chapter's hypothetical pool of \val{dn_records} records: recovery records are \val{dn_rec_pct} percent of records, \val{dn_stored_pct} percent of stored tokens, and \val{dn_target_pct} percent of selected target tokens, because a recovery record is longer (\val{dn_rec_len} stored tokens against \val{dn_dir_len}) and only its corrections are targets. Context per target is \val{dn_ctx_rec} for recovery against \val{dn_ctx_dir} for direct successes. Among the records a demonstration route admits, recovery is \val{dn_admitted_pct} percent of processed tokens (derived). Right, measured in the real corpus: resolved runs with a refused edit are \val{c_rec_pct} percent of resolved runs, \val{c_rec_chars_pct} percent of their characters, and \val{c_rec_ai_pct} percent of model-written characters; they are \val{c_len_ratio} times as long as direct successes on average (derived). ENGAGE: a team asks for ``\val{dn_rec_pct} percent recovery'' by record count. What did it actually buy? (Expected: nearly half the training compute and a different share of the gradient signal.) WARN: the right panel counts every model message as a target because the corpus masks by role; with assessed targets the shares would move. FLEX: cut if short; say the left panel's three numbers.}
\end{frame}

\begin{frame}{Mixed outcomes peak at mid difficulty}
  \begin{columns}[c]
    \begin{column}{0.56\textwidth}
      \centering\safeimg[width=\textwidth,height=4.6cm,keepaspectratio]{plateau.pdf}\\[0.02cm]
      {\footnotesize $1-(1-p)^{\val{gc_k}} - p^{\val{gc_k}}$ \quad (equation 3)}
    \end{column}
    \begin{column}{0.42\textwidth}
      \centering\safeimg[width=0.92\textwidth,height=3.8cm,keepaspectratio]{outcome_mix.pdf}\\[0.02cm]
      {\footnotesize Measured: outcomes per issue in the corpus}\\[0.1cm]
      \uncover<2->{\footnotesize At $p=\val{gc_lo}$: about \textbf{\val{gc_attempts_lo} attempts} per harvested success}
    \end{column}
  \end{columns}
  \mlsyscredit{Trajectory Curation, \S15.7.1, equation 3 and its figure. Right: nebius/SWE-agent-trajectories, CC BY 4.0; \val{c_tasks} issues; corpus\_stats.py.}
  \note{[2.5 min] [CORE] In a generate, select, train loop the current policy samples attempts, a check keeps the successes, and the next model trains on them: rejection-sampling fine-tuning. Left: with \val{gc_k} attempts at per-attempt pass rate $p$, the thin blue curve is the chance of at least one success, and the purple curve is the chance of a mixed group, at least one success and one failure. Mixed groups are what a preference pair or a group-relative objective needs. At \val{gc_lo_pct} and at \val{gc_hi_pct} percent the mixed fraction is \val{gc_mixed_lo}; at \val{gc_mid_pct} percent it is \val{gc_mixed_mid}. Step 2: at \val{gc_lo_pct} percent a fixture yields a success only \val{gc_yield_lo_pct} percent of the time, about \val{gc_attempts_lo} attempts per harvested success (derived). Right, measured: in the real corpus, \val{c_allfail_pct} percent of issues never succeed, \val{c_mixed_pct} percent are mixed, \val{c_allpass_pct} percent always succeed. ENGAGE: which fixtures should a contrast-hungry collector run more often? (Expected: the mid-difficulty ones; easy fixtures stay as cheap regression anchors, hard ones as coverage worth paying for.) WARN: all-success groups are still valid imitation material; equation 3 measures contrast and says nothing about usefulness. Groups with zero contrast and the gradient they lack belong to Lecture 17. Coverage and selection differ (Lecture 3): a loop that keeps only what it already solves narrows.}
\end{frame}

\begin{frame}{A published recipe filtered, deduplicated, and capped runs}
  \begin{columns}[T]
    \begin{column}{0.58\textwidth}
      \begin{logblock}[7.1]
        We collected a large number of successful\\ trajectories for easier problems but\\ significantly fewer for harder ones, which\\ required filtering and rebalancing ...\\[0.08cm]
        We first filtered the trajectories to\\ exclude those with clearly incorrect\\ actions, such as malformed LLM responses,\\ tool misuses, or linter-triggering errors.\\[0.08cm]
        ... retaining only unique action sequences,\\ and kept only the 30 shortest trajectories\\ for each problem.
      \end{logblock}
    \end{column}
    \begin{column}{0.40\textwidth}
      \small Applied to this corpus:\\[0.1cm]
      {\footnotesize Dropping every run with a linter refusal removes \textbf{\val{c_rec}} of \val{c_resolved} resolved runs (\val{c_rec_pct}\%), every recovery record.\\[0.12cm]
      Recovery runs are \textbf{\val{c_len_ratio}$\times$} as long as direct ones, so a shortest-first cap favors direct runs.}\\[0.15cm]
      \uncover<2->{\footnotesize\textbf{\textcolor{mlsysaccent}{Watch:}} distinct action sequences per fixture, held-out task families, an external anchor.}
    \end{column}
  \end{columns}
  \mlsyscredit{Nebius AI R\&D, ``Leveraging training and search for better software engineering agents,'' 2024-11-15 (excerpt, omissions marked). Counts: corpus\_stats.py on nebius/SWE-agent-trajectories, CC BY 4.0.}
  \note{[2.5 min] [OPTIONAL] A real recipe from the team that published the corpus, used in an expert-iteration loop. ``Many successes for easy problems, few for hard ones'' is self-selection narrowing, observed. Per-problem caps and unique action sequences are exactly the chapter's counter: they keep easy fixtures from flooding the corpus. The filter is the warning. We applied it to the same corpus as a counterfactual: removing every run with a linter-triggering error removes all \val{c_rec} resolved recovery records, \val{c_rec_pct} percent of the successes. Recovery runs average \val{c_rec_mean} characters against \val{c_direct_mean} for direct ones, \val{c_len_ratio} times as long, so keeping the shortest runs also favors the clean path, though with at most 30 per issue the recovery share falls only from \val{c_rec_pct} to \val{c_cap_rec_pct} percent. Step 2: the three signals that warn of narrowing are distinct normalized action sequences per fixture across rounds, single-attempt success rising while success within many attempts flattens, and held-out task families stalling while training fixtures climb. ENGAGE: was this filter wrong? (Expected: it depends on the objective. For a first imitation model, clean paths are reasonable; the cost is a model that never saw a refusal followed by a fix.) WARN: we applied their stated rules to the public corpus; their final training set may differ.}
\end{frame}

\begin{frame}{Vote: where does a success-only loop drift?}
  \mlsysactivity{Vote}{2 min}
  \centering
  {\normalsize Round after round, the policy trains on its own filtered successes. Training-fixture acceptance climbs. What is most likely?}\\[0.4cm]
  \begin{minipage}{0.86\textwidth}\normalsize
    \choice{A}{It keeps improving on held-out task families}\\[0.15cm]
    \answerat{2}{B}{It repeats strategies it already solves; held-out scores stall}\\[0.15cm]
    \choice{C}{It learns more recovery, since failures were filtered out}\\[0.15cm]
    \choice{D}{Nothing changes; the filter is unbiased}
  \end{minipage}\\[0.35cm]
  \uncover<2->{\small Bound it with per-task caps, an external anchor, and a held-out final test.}
  \mlsyscredit{Trajectory Curation, \S15.7.2; Shumailov et al., Nature 2024 (model collapse); Sculley et al., NeurIPS 2015 (hidden feedback loops).}
  \note{[2 min] [CORE] ENGAGE: vote, discuss with a neighbor for one minute, vote again. Step 2: B. Strategies that cleared earlier filters appear more often in training data, so the next policy samples them more, and the next filter keeps them again; diversity shrinks and the rare tail disappears first, the model collapse that Shumailov et al. measured. A is the hope; it holds only while held-out families keep rising. C has it backwards: filtering out failed steps removes recovery evidence. D ignores that the filter is part of the loop. Counter-measures: per-task retention caps and deduplication of identical action sequences; a permanent anchor of verified external demonstrations that self-play cannot overwrite; early stopping when diversity and held-out scores plateau; and a separate final test, because the benchmark consulted at every promotion becomes a development set. WARN: falling diversity alone can mean the agent cut useless retries; read it beside held-out transfer. A record goes stale when its tool schema, fixture, or grader changes, which can happen minutes after it was recorded. Land the segment aloud: a recorded mistake is context, and self-selected data narrows what it rewards.}
\end{frame}

% =============================================================================
% SECTION: COLLECTION
% =============================================================================
\mlsysnote{[0.5 min] [CORE] Last, the collector itself: how fast can it fill its buffers, and what must travel with every record it publishes?}
\mlsysdivider{How does the collector stay bounded, and what does the trainer read?}{Trajectory Curation, \S15.8 to \S15.10}
\section{Collection}

\begin{frame}{A verifier stall fills \val{st_store_gb} GB of staging in \val{st_tfull} s}
  \begin{columns}[c]
    \begin{column}{0.66\textwidth}
      \centering\mlsyssteps[width=\textwidth,height=4.7cm,keepaspectratio]{staging_step}{2}\\[0.02cm]
      {\footnotesize \val{st_rollouts} rollouts of \val{st_rollout_s} s; \val{st_slots} verifier slots; \val{st_ceiling_mb} MB ceiling per record}
    \end{column}
    \begin{column}{0.32\textwidth}
      \centering\safeimg[width=\textwidth,height=2.0cm,keepaspectratio]{atm-switch-forerunner.jpg}\\[0.03cm]
      {\footnotesize An ATM network switch. Credit flow control: send only while holding a free buffer slot.}\\[0.1cm]
      \uncover<2->{\footnotesize\textbf{Collection credit:} one per run, at the record ceiling, returned at durable disposition.}
    \end{column}
  \end{columns}
  \mlsyscredit{Trajectory Curation, Napkin Math 15.4 (hypothetical collector), figure 4; Kung et al., SIGCOMM 1994. Photo: Barcex, CC BY-SA 3.0, Wikimedia Commons (FORE Systems ForeRunnerLE 25).}
  \note{[3 min] [CORE] Collection is an asynchronous pipeline with staging between generation and verification. Nominally, \val{st_rollouts} rollouts of \val{st_rollout_s} s arrive at \val{st_arrival} records per second against \val{st_slots} slots that serve \val{st_capacity} per second: utilization \val{st_util}, which by Little's law (Lecture 5) is also the fraction of slots busy. Step 1: a defective fixture drives every test to its \val{st_timeout} s timeout. Capacity collapses to \val{st_burst} per second, so unassessed records grow by \val{st_net} per second, \val{st_growth} MB/s of reservations at the \val{st_ceiling_mb} MB ceiling. An empty \val{st_store_gb} GB store is full in \val{st_tfull_1} s, \val{st_tfull_min} minutes (derived). Step 2: stopping dispatch does not stop growth. Rollouts already running hold up to \val{st_inflight} GB, and \val{st_react_s} s of controller delay adds \val{st_react_gb} GB, so dispatch must stop at or below \val{st_mark_pct} percent. Right: credit-based flow control comes from network switches like this one: a sender forwards only while it holds a credit for a free downstream slot. A collection credit reserves the record ceiling through rollout, staging, and assessment, and returns only at a durable disposition or an acknowledged handoff. ENGAGE: why charge the ceiling? (Expected: the ceiling bounds the worst burst.) WARN: size staging for the service-time collapse, which average utilization hides. LINK: credit windows between cooperating agents are Lecture 18's.}
\end{frame}

\begin{frame}{A derived record inherits its split and permission}
  \centering
  \mlsyssteps[width=0.9\textwidth,height=5.15cm,keepaspectratio]{lineage_step}{2}\\[0.05cm]
  {\footnotesize Real corpus card: ``please respect the license of each specific repository'' (\val{c_repos} of them)}
  \mlsyscredit{Trajectory Curation, \S15.9, figure 5 (redrawn). Excerpt: nebius/SWE-agent-trajectories dataset card, License section, CC BY 4.0.}
  \note{[2.5 min] [CORE] Months later, someone asks which run produced a record, which grader version passed it, whether training on it is still permitted, and whether a redaction broke what justified the agent's next step. Step 1: a lineage envelope binds each record version to its collection identity, generation configuration, environment and capture, assessment, use and partition, and reverse references. A redaction never edits in place. trace\_r19\_01\_redacted\_v1 is a new record with a parent reference, the changed fields, the same split group, and the same use permission, and its assessment returns to pending until it is revalidated. The chapter's Example 15.1: if replacing a private path removes the clue the agent used, target suitability fails for that view. Step 2: when a source is withdrawn or a grader bug invalidates its checks, reverse references find every descendant. Bottom line: the real corpus carries a per-repository license condition, so use permission is a per-record field. ENGAGE: a manifest's MAC verifies. Does that prove the grader verdict is correct? (Expected: no. Integrity proves only unmodified bytes; truth and permission need their own evidence.) WARN: deleting a record from storage does not remove it from weights already trained; that needs retraining or audited unlearning.}
\end{frame}

\begin{frame}{Estimate: what does one emitted artifact cost?}
  \mlsysactivity{Estimate}{2 min}
  \begin{columns}[c]
    \begin{column}{0.56\textwidth}
      \centering\safeimg[width=\textwidth,height=4.2cm,keepaspectratio]{dispositions.pdf}
    \end{column}
    \begin{column}{0.42\textwidth}
      \small \val{co_attempts} attempts cost \$\val{co_cost} in all.\\[0.25cm]
      \wstep{2}{Emission yield: $\val{co_emitted}/\val{co_attempts} = \val{co_yield_pct}\%$}
      \wstep{2}{$\$\val{co_attempt_cost} \,/\, \val{co_yield} = $ \textbf{\textcolor{mlsysaccent}{\$\val{co_artifact}}} per artifact}
      \wstep{3}{At \val{co_dbl_pct}\% yield: \$\val{co_dbl_cost}}
    \end{column}
  \end{columns}
  \mlsyscredit{Trajectory Curation, Napkin Math 15.7 and equation 5 (hypothetical cohort). Values from images/calcs.py.}
  \note{[2 min] [OPTIONAL] ENGAGE: one minute. A cohort of \val{co_attempts} attempts costs \$\val{co_cost} for environments, generation, tools, checks, transforms, storage, and review. Every attempt appears once at the cutoff: \val{co_emitted} selected and emitted, \val{co_other} assessed but not selected, \val{co_inelig} ineligible for this use, \val{co_pending} pending, \val{co_unknown} unknown. Step 2: emission yield is \val{co_yield_pct} percent and the average attempt costs \$\val{co_attempt_cost}, so each emitted artifact costs \$\val{co_artifact} (derived): attempt cost divided by yield, the same form as cost per accepted task (Lecture 1). Step 3: doubling yield to \val{co_dbl_pct} percent halves the price to \$\val{co_dbl_cost}. ENGAGE: \val{co_not_pct} percent of attempts emit nothing. Where do you look first? (Expected: the \val{co_other} not selected. High redundancy means fewer attempts per task; ineligibility means split or permission defects caught too late; a pending backlog means a verification bottleneck.) WARN: emission yield is not task success; several non-selected runs solved their task but repeated a solution already kept. Cost says nothing about quality; only held-out evaluation decides that.}
\end{frame}

% =============================================================================
% CLOSE
% =============================================================================
\begin{frame}{What reaches the weights from \val{tri_fail} failures?}
  \mlsysactivity{Retrieval}{2 min}
  \centering
  \only<1|handout:0>{\safeimg[width=0.55\textwidth,height=2.0cm,keepaspectratio]{thumb_hook_step3.pdf}\\[0.15cm]
  {\normalsize Close your notes. Trace one record from a failure to a corpus manifest: name each decision and the evidence it needs.}}%
  \only<2->{\safeimg[width=0.95\textwidth,height=2.75cm,keepaspectratio]{synthesis.pdf}\\[0.15cm]
  \small
  \begin{tabular}{@{}P{3.6cm}P{9.6cm}@{}}\toprule
    Triage & \val{tri_system} repaired below the weights; \val{tri_residual} residual \\
    Split & task group assigned before collection \\
    Admission & R, U, P, O, T pass; tests out of the candidate's reach \\
    Targets & e3, e4, e6 selected; e2, e5 kept as context \\
    Lineage & redacted view re-admitted; split and permission inherited \\\bottomrule
  \end{tabular}}
  \mlsyscredit{Trajectory Curation, \S15.10 (case study) and Napkin Math 15.1.}
  \note{[3 min] [CORE] Back to the hook. Step 1: retrieval, 90 seconds. Step 2: the \val{tri_fail} failures of the hook. Probes route \val{tri_system} to a schema fix and a larger sandbox pool, below the weights, for every model. The \val{tri_residual} that survive define a learning objective: here, recovering after a harness refusal. Before collection, the task group pool\_repair\_01 is assigned to training and its final-test relatives are protected. The recovery run collection\_run\_r19\_01 produced ws\_r20, which satisfies repair\_accept\_v1. Its record passes the five checks for the demonstration route, with tests out of the candidate's reach. The target manifest selects e3, e4, and e6 and keeps e2 and e5 as context. Redacting a path creates trace\_r19\_01\_redacted\_v1, which inherits the split group and permission and is re-admitted on its own evidence. The manifest repair\_demonstrations\_v1 binds all of it for the trainer. ENGAGE: can the same record become a preference pair? (Expected: no. e2 and e4 follow different prefixes, so they are not alternative continuations of one context.) WARN: none of this proves the trained model is better. Only held-out evaluation on protected task groups does. FLEX: [CORE].}
\end{frame}

\begin{frame}{Takeaways}
  \thumbline{thumb_ladder_step3.pdf}{Probes send most failures to context, schemas, or the runtime; training gets the residual (\val{tri_residual} of \val{tri_fail} in the hook).}\\[0.12cm]
  \thumbline{thumb_groups_step2.pdf}{A random split of runs leaks most tasks (\val{c_straddle_pct}\% of real issues); split whole task groups before collecting.}\\[0.12cm]
  \thumbline{thumb_cascade_step3.pdf}{Admission is five checks per route; order them by cost per rejection, never by dropping one.}\\[0.12cm]
  \thumbline{thumb_typology_step3.pdf}{A recorded mistake stays as context; self-selected data narrows toward what it already solves.}\\[0.12cm]
  \thumbline{thumb_staging_step2.pdf}{Size staging for a verifier stall with credits, and give every derived record its own lineage and decision.}
  \note{[1 min] [CORE] One line per segment. Ask students which slide taught each line. If short on time, read the first and the fourth.}
\end{frame}

\begin{frame}{Open questions}
  \begin{columns}[c]
    \begin{column}{0.65\textwidth}
      \small
      \begin{itemize}\setlength{\itemsep}{0.12cm}
        \item Public corpora mask every model turn. Who should pay to assess targets step by step: the publisher or each team that trains?
        \item When a source repository withdraws permission, what does a team owe for models already trained on it?
        \item Can a self-improvement loop ever stop needing an external anchor of human demonstrations?
      \end{itemize}
    \end{column}
    \begin{column}{0.33\textwidth}
      \centering\mlsysphoto[width=\textwidth,height=3.4cm,keepaspectratio]{datacenter-server-racks.jpg}\\[0.05cm]
      {\footnotesize Server racks: where rollouts, checks, and training all run}
    \end{column}
  \end{columns}
  \mlsyscredit{Photo: Carl Lender, CC BY 2.0, Wikimedia Commons.}
  \note{[1 min] [OPTIONAL] No single answer. Pick one for discussion or assign one as a short written response. The first follows from the Diagnose slide: assessed targets cost review time, and role masks are free. The second follows from \S15.9.4: deleting a record leaves completed weight updates intact. The third follows from \S15.7.2 and model collapse.}
\end{frame}

\begin{frame}{Read and explore}
  \small
  \begin{itemize}\setlength{\itemsep}{0.03cm}
    \item \textbf{Read:} Trajectory Curation, \S15.2 to \S15.10; Checkpoints 15.1 and 15.2; Napkin Math 15.1 to 15.4 by hand; exercises 1, 2, and 4.
    \item Saltzer, Reed, and Clark, ``End-to-End Arguments in System Design,'' 1984.
    \item Hellerstein and Stonebraker, ``Predicate Migration,'' SIGMOD 1993.
    \item Ross, Gordon, and Bagnell, ``A Reduction of Imitation Learning \ldots,'' AISTATS 2011.
    \item Gebru et al., ``Datasheets for Datasets,'' CACM 2021.
    \item Shumailov et al., ``AI Models Collapse When Trained on Recursively Generated Data,'' Nature 2024.
  \end{itemize}
  \note{[0.5 min] [CORE] The chapter first, with its checkpoints and notebooks. Saltzer, Reed, and Clark give the end-to-end argument for placing a function where it can be done correctly. Hellerstein and Stonebraker order expensive predicates by rank, the cascade rule. Ross et al. explain why clean demonstrations leave a policy unprepared for its own mistakes. Gebru et al. is the documentation discipline a corpus manifest automates. Shumailov et al. measure collapse under recursive training. FLEX: assign images/data/corpus\_stats.py as a lab: download the corpus, rerun the split, and change the cap.}
\end{frame}

\begin{frame}[c]
  \begin{columns}[c]
    \begin{column}{0.34\textwidth}
      \centering\safeimg[width=\textwidth,height=2.4cm,keepaspectratio]{thumb_typology_step3.pdf}
    \end{column}
    \begin{column}{0.62\textwidth}
      {\small\textcolor{lightgray}{NEXT LECTURE \textbar{} TRAJECTORY FINE-TUNING}}\\[0.5cm]
      {\LARGE\bfseries The record is admitted. Which of its tokens should the loss touch?}
    \end{column}
  \end{columns}
  \note{[0.5 min] [CORE] Lecture 16, Trajectory Fine-Tuning, opens with the record this lecture published to repair\_demonstrations\_v1. Observations are most of a stored trajectory, and an unmasked loss trains the model to write its own tool results. Next time: serializing the exact prefix the agent saw, masking the loss to the selected targets, and deciding whether a lower loss earns a release.}
\end{frame}

\end{document}
